\chapter{\rm\bfseries Introduction}
\label{ch:introduction}

\section{Motivation}

Fabric handling automation is a long-standing open problem in industrial
robotics, with applications ranging from apparel manufacturing and laundry
processing to packaging of textile components. Cloth is uniquely difficult to
grasp robustly: it deforms continuously under contact, exhibits no fixed
graspable feature, and has a near-zero bending stiffness that allows it to slip
out of a poorly placed gripper. Conventional rigid-jaw grippers, designed for
sturdy objects with well-defined geometry, perform poorly on cloth without
elaborate tactile sensing or human-tuned approach trajectories.

A class of compliant grippers based on the finray effect---a passive
bending mechanism inspired by fish-fin anatomy---has been proposed as a
solution. Finray fingers passively wrap around the contact surface, increasing
contact area and gripping force without requiring active force control. This
compliance, however, comes at the cost of force-modelling difficulty: the
finger geometry deforms under load, the contact point shifts, and tilt-induced
asymmetries arise that have no analogue in the rigid-jaw case.

The central question of this thesis is: \emph{to what extent does the finray's
compliance trade single-axis force precision for multi-axis operational
robustness?} Specifically, how does the finray's contact force respond to
deviations in the two design degrees of freedom most likely to be uncontrolled
in real deployment: press depth and gripper tilt angle?

\section{Research Question and Contributions}

This thesis answers the following research questions through a controlled
experimental benchmark on a Franka FR3 collaborative robot:
\begin{enumerate}
  \item How does stationary contact force vary with press depth for each of
        five candidate gripper geometries?
  \item For each geometry, what is the maximum tilt angle (about each
        orthogonal axis) at which a reliable grasp is still produced?
  \item Is the dependence of contact force on tilt itself depth-dependent?
  \item Does the finray gripper exhibit a direction-dependent (handed)
        response under tilt about its geometrically asymmetric axis?
\end{enumerate}

\noindent The contributions of this work are:
\begin{itemize}
  \item A reproducible benchmark protocol for cloth-grasp force characterisation
        on a Franka FR3, with full source code published alongside this thesis.
  \item A two-dimensional depth--tilt experimental matrix sampling five
        commercially relevant gripper geometries.
  \item Empirical evidence of a depth-dependent tilt robustness transition
        in finray grippers, with quantitative characterisation of the
        transition depth.
  \item A documented calibration procedure for per-variant tool-centre-point
        (TCP) offset using differential force touch-off against a rigid
        reference.
  \item A geometric verification framework that disentangles control-frame
        positioning error from physical gripper compliance under load.
\end{itemize}

\section{Thesis Organisation}

Chapter~\ref{ch:background} reviews the literature on cloth grasping,
compliant grippers, and force-controlled approach strategies.
Chapter~\ref{ch:methods} describes the experimental hardware, the kinematic
calibration procedure, and the data-acquisition pipeline.
Chapter~\ref{ch:experiments} presents the depth-sweep characterisation (Experiment~1)
across all five gripper variants at zero tilt.
Chapter~\ref{ch:experiments} presents the tilt-sweep characterisation (Experiment~2)
at boundary and safe depths.
Chapter~\ref{ch:discussion} synthesises the findings into a model of
depth-dependent tilt robustness.
Chapter~\ref{ch:conclusion} concludes and identifies directions for future work.
